Entre les dues làmines de la figura, separades una distància $d = 3,0\ \mathrm{m}$, tenim un camp elèctric uniforme de $1,5\times 10^{3}\ \mathrm{N\,C^{-1}}$.

\figura[0.32]{fig1.png}

En el centre de l'espai limitat per les dues làmines posem una llentia metàl·lica carregada, penjada d'un fil. Tenint en compte que la longitud del fil és de 1,5~m, que la càrrega de la llentia és de $Q = -5,0\times 10^{-5}\ \mathrm{C}$ i que té una massa $m = 12\ \mathrm{g}$:

\begin{apartat}{1}
Representeu les forces que actuen sobre la llentia en el punt d'equilibri i calculeu l'angle que forma el fil amb la vertical en l'equilibri.
\end{apartat}

\begin{apartat}{1}
Calculeu la diferència de potencial entre la posició d'equilibri i la posició vertical.
\end{apartat}
